% =====================================================================
%  Professioni 1 — I Dev Days: il metodo di AIV e i Dev Projects
%  Fonte: trascrizione dell'incontro di apertura dei Dev Days (AIV1):
%  un incontro online su Zoom di un'ora, il lunedì mattina della prima
%  settimana del primo anno, tenuto da Daniele Fusetto (direttore del
%  corso di narrative design e brand advisor dell'Accademia Italiana
%  Videogiochi) con Andrea Morini (direttore della didattica). Circa 31
%  partecipanti collegati su un centinaio di iscritti al primo anno;
%  slide condivise, un video di gioco senza audio, domande in chat.
%  Pacchetti richiesti nel preambolo del libro:
%    \usepackage[utf8]{inputenc}, \usepackage{booktabs},
%    \usepackage{tikz} con
%    \usetikzlibrary{arrows.meta,positioning,calc,fit}
%
%  Collocazione nel libro: parte "Alcune Professioni nel Mondo
%  Videoludico", dopo AIV0_mindset_e_best_practice (3.2).
%  Tutte le etichette di questo capitolo hanno il prefisso "p1:".
%  Il capitolo AIV0 (prefisso "p0:") rimanda a questo; questo rimanda ai
%  capitoli delle singole professioni (AIV2-AIV6).
%  La figura è una ricostruzione schematica di quanto il relatore spiega
%  con le slide, non una copia delle slide.
% =====================================================================

\chapter[I Dev Days e il metodo di AIV]{I Dev Days: il metodo di AIV,
la collaborazione fra discipline e i Dev Projects}
\label{p1:cap:devdays}

Ogni anno accademico, all'Accademia Italiana Videogiochi (\textbf{AIV}),
comincia con una settimana di incontri chiamati \textbf{Dev Days}. Lo scopo
è duplice: \textbf{impostare il lavoro} dell'inizio dell'anno e
\textbf{presentare a tutti gli studenti del primo anno che cosa faranno i
colleghi degli altri corsi}, perché, come si sentirà ripetere molte volte,
``la gaming industry è fortemente collaborativa'': i giochi si fanno in team,
e per lavorare bene con un game designer, un narrative designer, un
animatore, un programmatore o un 3D artist bisogna sapere che cosa fanno.
Questo capitolo raccoglie l'\textbf{incontro di apertura}; i capitoli
successivi raccolgono le presentazioni dei singoli corsi, cioè delle singole
professioni: game e level design (capitolo~\ref{p2:cap:gamedesign}),
narrative design (capitolo~\ref{p3:cap:narrative}), character animation
(capitolo~\ref{p4:cap:animation}), programmazione
(capitolo~\ref{p5:cap:programming}) e 3D art (capitolo~\ref{p6:cap:3d}).

A condurre l'incontro è \textbf{Daniele Fusetto}, direttore del corso di
narrative design e \emph{brand advisor} di AIV, che insegna in AIV dal 2019 e
in generale dal 2016; con lui c'è \textbf{Andrea Morini}, \textbf{direttore
della didattica}, che interviene a più riprese. L'incontro si tiene online,
su Zoom, il lunedì mattina, e dura \textbf{un'ora} (``abbiamo un'ora di tempo,
purtroppo''): non si aspetta molto chi è in ritardo, anche perché la
registrazione verrà caricata sulla piattaforma dell'accademia
(\emph{l'hub}). L'incontro non è facoltativo: ``i Dev Days non sono degli
incontri facoltativi'', bisogna almeno vederli in registrazione, ``altrimenti
si rischia di rimanere indietro fin dall'inizio''. Negli anni in cui
Fusetto e Morini hanno insegnato, ``sottovalutare una cosa del genere è di
solito il primo sintomo che qualcosa non va''.

Due avvertenze di lettura. La prima: la trascrizione è automatica e storpia
quasi sempre il nome dell'accademia (``Hive'', ``IVE'', ``Ive'', ``AIB'',
``AB'', ``IV'') e diversi titoli e termini; il capitolo usa sempre
``AIV'' e segnala in nota le ricostruzioni incerte. La seconda: molte
affermazioni riguardano l'organizzazione di AIV (calendari, regole,
certificati, opzioni di iscrizione) e valgono per l'anno a cui l'incontro si
riferisce; sono riportate come le dicono i relatori.

La data, dedotta, è \textbf{lunedì 5 ottobre 2026}\footnote{La data non è
detta. Si ricava da tre indizi: le lezioni dei corsi cominciano
``ufficialmente'' il 19 ottobre, quelle di narrative ``questo venerdì''; nel
capitolo sul narrative design, il giorno dopo, il relatore fissa la prima
lezione di narrative a ``venerdì 9 ottobre''; e nel 2026 il 9 ottobre è
venerdì. È una deduzione.}.

Il percorso del capitolo:
\begin{enumerate}
  \item l'avvio e la politica della webcam (sezione~\ref{p1:sec:avvio});
  \item che cosa sono i Dev Days e il calendario della settimana
        (sezione~\ref{p1:sec:devdays});
  \item il \textbf{metodo di AIV}: che cosa non si tollera e quattro cose da
        fare (sezione~\ref{p1:sec:metodo});
  \item \textbf{best practice} e \textbf{trappole}
        (sezioni~\ref{p1:sec:best} e~\ref{p1:sec:trappole});
  \item gli \textbf{obiettivi} dei corsi: portfolio, lavoro di gruppo,
        competenze trasversali (sezione~\ref{p1:sec:obiettivi});
  \item l'attestato di partecipazione e il \textbf{certificato di
        eccellenza} (sezione~\ref{p1:sec:eccellenza});
  \item i \textbf{Dev Projects} del secondo anno e la pipeline di sviluppo,
        con il caso di \emph{Beat of Rebellion}
        (sezioni~\ref{p1:sec:devprojects}--\ref{p1:sec:beat});
  \item perché parlarne il primo giorno, i falsi miti e le domande finali
        (sezioni~\ref{p1:sec:perche}--\ref{p1:sec:domande}).
\end{enumerate}

% ---------------------------------------------------------------------
\section{L'avvio}
\label{p1:sec:avvio}

Si comincia con qualche minuto di attesa: alcune persone sono in ritardo
``per problemi di hub'' o perché non hanno ancora fatto l'accesso. Il
relatore disabilita la sala d'attesa di Zoom e raccomanda di tenere
d'occhio il \textbf{Discord di AIV} e la propria \textbf{mail}. Poi
condivide le slide e comincia con la parte ``molto introduttiva'', mentre
arrivano gli ultimi.

\subsection{La webcam}

``Possibilmente accendete la webcam'': il relatore vuole vedere facce sveglie
``e anche le vostre camerette''. ``Mamma mia, è lunedì mattina per tutti'':
non c'è una regola che imponga di presentarsi a lezione ``belli e col volto
riposato'', ma la regola è che \textbf{ci si fa vedere}, anche quando si
segue da remoto. Non solo per buona educazione: i docenti devono vedere le
facce, le reazioni, e vogliono conoscere gli studenti anche fisicamente. Già
al tempo del COVID, ricorda, riconoscere dal vivo le persone viste online o
con la mascherina era stato ``un dramma''.

% ---------------------------------------------------------------------
\section{Che cosa sono i Dev Days}
\label{p1:sec:devdays}

L'agenda della mattina: che cosa sono i Dev Days, qual è il \textbf{metodo di
AIV}, e i \textbf{Dev Projects}, ``una cosa che voi vedrete al secondo
anno'', ma di cui i relatori vogliono dare un assaggio fin da oggi ``per
farvi capire che qua non si scherza''\footnote{La trascrizione dice
``farvi la capatanza'' e poi ``capatanta'': il senso è ``farvi fare una
capatina'', dare un assaggio. Il nome dei progetti del secondo anno
oscilla nella trascrizione fra ``Game Dev'', ``Dev Project'' e ``The
Projects''; il capitolo usa ``Dev Projects'' per i progetti e ``Game Dev''
per il percorso del secondo anno in cui si svolgono.}.

\subsection{Un'industria che lavora in team}

I Dev Days sono incontri dedicati, all'inizio del primo anno, a presentare
ciò che faranno gli altri colleghi e a impostare il lavoro. Il motivo è che
l'industria è fortemente collaborativa: \textbf{si sviluppano giochi in
team}, ``dalle 20 persone alle 400 persone''. Ci sono anche team più
piccoli, ma spesso un team di tre, quattro, cinque persone lavora molto con
\textbf{freelancer}. Esempio: si dice che \emph{Clair Obscur: Expedition 33}
sia stato fatto da trenta persone; ``non è vero, se vedete i credit, sono
400''\footnote{La cifra di 400 è del relatore. Il nucleo dello studio
Sandfall Interactive era di circa trenta persone; i titoli di coda includono
però numerosi collaboratori esterni (musica, doppiaggio, animazione,
outsourcing, localizzazione), che è appunto il punto del relatore.}. Fare
un gioco richiede un lavoro di team.

Al primo anno ognuno fa la propria disciplina: chi fa programmazione fa
programmazione, chi fa 3D fa 3D, e così per game design, character
animation, narrative design. Poi, all'inizio del secondo anno, ci si
incontrerà con tutti gli altri ``per fare insieme un videogioco''. Ma fin da
oggi bisogna capire che cosa faranno i colleghi e quali sono le altre
professioni insegnate in AIV, ``per meglio collaborare con gli altri''.

\subsection{Il calendario della settimana}

Per questo la settimana prevede \textbf{cinque incontri}, uno per ciascun
corso di AIV. Vanno seguiti \textbf{tutti}, e non sono solo per gli studenti
del corso presentato: ``sono per tutti gli studenti''. La presentazione di
game design spiega che cosa faranno i game designer durante l'anno, e la
devono seguire tutti, così che l'anno dopo, collaborando con loro, si sappia
più o meno che cosa fanno di solito. Lo stesso per narrative design,
character animation, programmazione e 3D art. L'obiettivo non è spiegare agli
studenti di quel corso che cos'è il corso (lo vedranno durante l'anno), ma
spiegare \textbf{a tutti} che cosa fanno le altre discipline. Tanto che
``probabilmente è meglio se, per esempio, la presentazione del corso di
programming la seguono più attentamente quelli che non faranno
programming''.

In più:
\begin{itemize}
  \item un incontro aggiunto quest'anno, il \textbf{giovedì mattina}, di
        un'ora e mezza, con Fusetto, anch'esso obbligatorio, sulla
        \textbf{comprensione del testo} (capitolo~\ref{pt:cap:comprensione}): ``mi spiace dirlo, ma negli ultimi
        anni ci siamo accorti che è difficile arrivare all'inizio di AIV con
        delle capacità di comprensione del testo ottime''. Non si parla di
        italiano di base né di una disciplina, ma di \textbf{comprendere
        ciò che le persone dicono e vogliono}, rispondere correttamente alle
        richieste dei docenti, evitare incomprensioni; si vedranno anche
        trucchi per comprendere meglio e per segnarsi le cose in modo da
        \textbf{saper porre domande}, ``perché è importantissimo porre
        domande'';
  \item un incontro il \textbf{martedì pomeriggio alle 14:30}, solo per chi
        si è iscritto entro il 31 luglio con la \textbf{Gamer Promo} o la
        \textbf{Dev Promo}: un incontro di \emph{mindset} e \emph{best
        practice} sulle pratiche migliori per imparare e apprendere
        durante l'anno (è il capitolo~\ref{p0:cap:mindset});
  \item per gli stessi studenti, la possibilità di una
        \textbf{mentorship} di un'ora con un tutor o un docente, ``per
        partire al meglio''. ``Sfruttatele, queste ore, perché l'anno scorso
        nessuno l'ha usata e poi dopo saltano fuori i problemi.'' Chi non
        ricorda se ha la promo può chiedere in segreteria.
\end{itemize}

% ---------------------------------------------------------------------
\section{Il metodo di AIV}
\label{p1:sec:metodo}

\subsection{In AIV si boccia}

``In AIV si boccia.'' Non è un'accademia ``dove vieni, paghi e passi
l'anno'': ci sono dei criteri. Il criterio principale è \textbf{non
consegnare le attività}, le cosiddette \textbf{task}, richieste durante
l'anno, perché la valutazione si fa su quelle: se non vengono consegnate, i
docenti non possono valutare.

Ma contano anche le \textbf{soft skill}, ``e se avete frequentato il nostro
open day lo sapete'': per passare l'anno bisogna
\begin{itemize}
  \item \textbf{interagire}: non isolarsi; chi è un po' introverso cerchi di
        uscire, di chiedere, di porre domande;
  \item \textbf{rispettare gli altri};
  \item essere \textbf{proattivi}: se viene in mente oggi che si deve o si
        può fare una cosa, non la si rimanda a domani; se si ha una domanda
        per la prossima lezione e si sa che il docente è raggiungibile per
        mail o per messaggio, non si aspetta la lezione.
\end{itemize}

\subsection{Tolleranza zero: bullismo, pregiudizi, cheating}

Su tre cose non c'è tolleranza.
\begin{description}
  \item[Bullismo.] ``Self-explanatory'', in tutte le sue forme: qualunque
        prevaricazione, anche implicita, ``non esiste''. Massima
        collaborazione e massima gentilezza, sempre.
  \item[Pregiudizi.] Si è in tanti e molto diversi, ognuno con le sue
        opinioni; ``a meno che voi non abbiate problematiche di tipo,
        diciamo, criminale, ma non credo'', bisogna partire dal presupposto
        che la persona che si ha davanti ha una sua storia, mettersi nei
        panni degli altri e \textbf{sospendere il giudizio}.
  \item[Cheating.] Fino a qualche anno fa voleva dire copiare dai compagni,
        avvantaggiarsi, manipolare la classe o i compagni per ottenere
        vantaggi non legati alla performance (``non faccio bene le attività,
        ma sono bravo a manipolare l'ambiente in modo da passare''). Negli
        ultimi due anni si è aggiunto \textbf{l'uso dell'IA}: è cheating
        quando, in un corso, viene detto di non usarla in un certo modo e la
        si usa così. Ogni corso ha il suo approccio (programmazione diverso
        da narrative, da game design, da character animation) e verrà detto
        che cosa si può fare e che cosa no. In narrative design, per
        esempio, è \textbf{severamente vietato scrivere con l'IA} per tutti
        e tre gli anni, ``vietatissimo'', fino a una lezione in cui il
        relatore spiegherà ``qual è l'unico modo, a mio parere'' per farsi
        aiutare dall'IA nell'atto della scrittura.
\end{description}

Ma cheating è anche copiare dai compagni: ``è successo, ce ne accorgiamo''.
I docenti vedono i compiti, notano forme usate nelle attività o
``spennellature, scolpiture nel modello 3D'' troppo simili. ``D'altronde, se
molti critici hanno capito che Palworld aveva copiato da Pokémon anche i
modelli 3D, si può fare''\footnote{È l'opinione del relatore, riportata
com'è. Le somiglianze fra i modelli di \emph{Palworld} e quelli dei Pokémon
sono state sostenute da utenti e commentatori dopo l'uscita del gioco
(2024), ma non sono mai state accertate; la causa intentata da Nintendo e
The Pokémon Company contro lo sviluppatore riguarda la violazione di
brevetti, non la copia dei modelli.}. Nel cheating c'è sempre una vittima,
la persona sfruttata o manipolata; ma soprattutto chi bara \textbf{fa un
danno a se stesso}: ``state pagando per imparare. Imparare voi.'' Chi usa
scorciatoie per passare l'anno e poi si ritrova senza conoscenze ha perso
un'occasione.

\subsection{Inquiry-based learning}

Poi ``quattro cose che dovete fare''. La prima: tutti i corsi di AIV sono
\textbf{inquiry-based}. ``Non esistono domande stupide, soprattutto il primo
anno.'' Molti studenti scrivono ``ho una domanda stupida'': no, non ce ne
sono. Arriva proprio in quel momento una domanda in chat (``scusi, ma i
corsi che ha fatto vedere prima sono online?''), e il relatore la usa come
esempio: ``anche qua, ragazzi, è importantissimo, \textbf{specificate bene
nelle domande}''. Si intende la \emph{lezione}, cioè il singolo incontro, o
il \emph{corso}, cioè tutto il corso? Gli incontri di questa settimana sono
tutti da remoto. Sara ha fatto una domanda ``per niente stupida'', perché
non era chiaro, ``e ha imparato anche qualcosa''\footnote{Nel capitolo sul
narrative design (sezione~\ref{p3:sec:skill}) il relatore torna su questo
episodio per spiegare perché tiene tanto alla precisione dei termini.}.
``Meno domande fate, meno imparate.''

\subsection{Un ambiente simulativo del mondo del lavoro}

La seconda: AIV è un \textbf{ambiente simulativo del mondo del lavoro}.
Tranne game programming, che ha test ed ``esamini'' (e i docenti di quel corso
spiegheranno perché), \textbf{non ci sono sessioni d'esame}, ma valutazioni
verbali tramite \textbf{task}. Durante molte lezioni, non tutte, il docente
dirà: ``voglio che mi consegniate questa task entro questa data''. Se la si
manca, la task è mancata e conta per il non passaggio dell'anno. Le task non
hanno tutte la stessa durata: in alcuni corsi sono molto corpose e
richiedono mesi; in game design e narrative sono più equilibrate, con task
da una o due settimane e altre da mesi. In sintesi: task settimanali,
mensili, trimestrali o annuali.

Bisogna entrare a poco a poco nell'ottica di un \textbf{approccio
professionale} al corso: ``prendetelo un po' come un lavoro''. Ogni giorno
si deve produrre, studiare, apprendere, confrontarsi, fare call, come in
qualsiasi lavoro, non solo nei videogiochi. La puntualità è importante;
tutti gli aspetti professionali ``in senso lato'' non sono semplicemente
graditi, ``sono mandatori''. ``Molto di più di quando i genitori ci dicevano,
alle superiori, devi pensare alla scuola come a un lavoro: questo è molto più
veritiero.'' E come nel lavoro si fanno meeting e incontri, ma si lavora
anche molto da soli: non bisogna sottovalutare il peso delle attività.

Le task sono \textbf{opportunità} per migliorarsi, ottenere più feedback e
\textbf{farsi notare dai docenti}, perché il 99\% dei docenti di AIV sono
lavoratori del settore: Fusetto stesso lavora nell'industria, sta facendo dei
giochi in questo periodo, e alcuni suoi studenti lavorano con lui. ``Non è un
modo per essere assunti'', ma permette di farsi notare per le proprie
capacità.

\subsection{Fail to learn}

La terza: \textbf{fallire}. Una delle cose più importanti da fare al primo
anno è fallire e imparare: non \emph{fail for}, ma \emph{fail to
learn}\footnote{Così nella trascrizione; il senso è ``fallire per
imparare''.}. Se si sbaglia, si ricevono feedback e si trovano soluzioni, e
piano piano si impara a trovare soluzioni ai propri fallimenti; si noterà
``un aumento notevole'' delle proprie capacità. Il consiglio pratico:
prendere una task di settembre, metterla ``in una giara virtuale'' (un
drive), lasciarla lì e a gennaio confrontarla con una task di gennaio. Se le
differenze non si notano, è perché non si è fallito e non si è imparato: ``è
una campanella d'allarme''.

\subsection{Mettere in pratica, ogni giorno}

La quarta: i corsi di AIV richiedono di \textbf{mettere in pratica} ciò che
si apprende. Perché solo \textbf{nove ore di lezione a settimana}, e non
lezioni tutte le mattine come in altre accademie? Perché bisogna avere il
tempo di mettere in pratica teoria e pratica viste a lezione: ``se io non vi
do almeno 20 ore alla settimana per mettere in pratica quello che vi insegno,
è difficile''. Il relatore porta il proprio esempio\footnote{Poco prima il relatore dice ``come diceva Andrea'': nella trascrizione i cambi di voce fra i due relatori non sono segnati, e l'attribuzione di alcuni passaggi è incerta. Qui, dal contesto (``io so programmare un pochettino'', ripetuto anche nel capitolo~\ref{p3:cap:narrative}), parla con ogni probabilità Fusetto.}: pur considerandosi un artista usa spesso Unreal Engine, e nel lavoro si acquisiscono capacità trasversali; sa
programmare ``un pochettino'', ma non ci riuscirebbe se non si mettesse
tutti i giorni quindici minuti a programmare\footnote{La trascrizione dice
``io non saprei programmare se io tutti i giorni mi mettessi lì 15 minuti a
programmare'': manca evidentemente un ``non''.}. Bisogna impegnarsi ogni
giorno nella propria disciplina. Altrimenti ci si sveglia all'ultimo, si
imparano le cose al momento del bisogno (magari anche bene), ma nasce una
sensazione di \emph{burnout}, ``e poi il burnout lo prendete''.

Quindi: mettersi da parte del tempo ogni settimana per applicare ciò che si
vede a lezione, farsi domande, fare domande ai docenti, e
\textbf{impegnarsi al massimo sulle task}. Con alcune regole:
\begin{itemize}
  \item se una task sembra \textbf{semplice}, la si fa \textbf{due volte}:
        forse si sta sopravvalutando se stessi. Se è davvero semplice, si
        chiede al docente se c'è la possibilità di avere qualcosa in più
        (``a me sembrava semplice, secondo te l'ho fatta bene?''), e magari
        arriva qualcosa di extra;
  \item se sembra \textbf{troppo difficile}, la si tenta comunque, perché
        una delle basi dell'apprendimento è l'\textbf{alternanza di
        difficoltà}: alternare task semplici e difficili fa apprendere
        meglio che fare solo le une o solo le altre;
  \item se la si finisce \textbf{in anticipo}, la si approfondisce e la si
        rifà. ``È noioso? Sì. È lungo? Sì. È pesante? Sì. Però è il
        migliore.''
\end{itemize}

\subsection{Consegnare in anticipo per avere un feedback?}

Una domanda in chat: si può consegnare una task in anticipo per ricevere un
feedback e correggere gli errori? Sì, risponde il relatore, ma con due
precisazioni. Di solito i docenti danno tutti i feedback contemporaneamente,
``perché ho bisogno di vedere anche il confronto con la classe''; si può
però mandare una task e fare \textbf{domande specifiche}. E bisogna
aspettarsi che i docenti non rispondano direttamente: a lui capita spesso
che gli chiedano ``Daniele, secondo te questa cosa l'ho spiegata troppo?
l'ho fatta troppo così?'', e lui risponde ``secondo te l'hai fatta
troppo?''. ``Non frustratevi per questo'': deve insegnare in qualche modo ad
andare avanti da soli. Non è detto quindi che si ottenga una risposta, ma si
avranno indicazioni, e chi consegna uno o due giorni prima può ricevere
feedback aggiuntivi. Una condizione: se si consegna in anticipo, deve
essere perché si è dato il massimo, non perché si è ``tirato via'' per
consegnare prima. E non è banale: nella sua carriera, ``ma anche
attualmente'', quasi nessuno riesce a consegnare in anticipo.

% ---------------------------------------------------------------------
\section{Best practice}
\label{p1:sec:best}

Le \textbf{best practice}, ``soprattutto per chi segue in remoto''.

\paragraph{Camera sempre accesa.} Vale per chi segue da remoto, ma anche per
chi segue in presenza e capita di seguire una lezione online. E capiterà:
``in Italia abbiamo una maledizione, si chiamano scioperi''. Spesso il
venerdì ci sono scioperi di mezzi, treni, autostrade, a Roma e a Milano
(``l'anno scorso Milano è stato un delirio''); per non perdere la lezione,
se c'è sciopero e il docente non può venire in presenza, la lezione del
venerdì si trasforma in lezione da remoto. In quei casi la camera va tenuta
accesa: al docente serve per vedere i volti (``vedo il volto di Mattia, di
Gabriele, di Emanuele, di Davide, vedo come sono messi, se stanno capendo
oppure no''), e allo studente è ``una scusa in meno per distrarvi'', perché
``la distrazione è sempre dietro l'angolo''.

\paragraph{Prendere appunti.} Cartacei o digitali, non importa. Non bisogna
fare come i camerieri, che pensano di prendere gli ordini a mente e
puntualmente se li dimenticano; il fenomeno ha anche un nome, l'\textbf{effetto
Zeigarnik}, ``se volete potete andarlo a cercare''\footnote{La trascrizione
dice ``Zygarnik Effect''. L'effetto è ripreso e discusso nel
capitolo~\ref{p0:cap:mindset} (sezione~\ref{p0:sec:memoria}).}.

\paragraph{Interagire con i docenti.} Fare domande durante le lezioni.

\paragraph{Rivedere le registrazioni.} Soprattutto per gli
\textbf{studenti lavoratori}. Ogni settimana ci sono nove ore di lezione; chi
lavora e sa che ne perderà tre a settimana ha il compito ``tassativo'' di
recuperarle nel weekend. Altrimenti la settimana dopo diventano sei, poi
nove, ``e improvvisamente quelle tre ore diventano 18''. ``Io ho studenti che
arrivano a giugno e mi dicono: prof, devo ancora recuperare lezioni di
febbraio. E io non rido'', perché le lezioni sono spesso consequenziali. Il
sistema di registrazione dà pochi problemi: entro 24 ore circa le
registrazioni sono disponibili (``scongiurando bug e problematiche di
cloud''). ``Guardatevele, riguardatevele, mettete il 2x piuttosto.''

\subsection{Quanto restano le registrazioni, e a che cosa servono}

A una domanda in chat uno dei relatori risponde che le registrazioni restano
attive almeno fino all'inizio dell'anno successivo, verso fine settembre,
quando si cominciano a staccare le vecchie per fare spazio al nuovo anno.
Si possono quindi recuperare anche d'estate: chi, per motivi di lavoro, non
ha consegnato abbastanza e ha bisogno di tempo per recuperare, lo ottiene
durante l'estate, con l'accesso alle lezioni.

Le registrazioni, aggiunge Fusetto (richiamando un'osservazione di Morini), servono anche ad altro: ci
saranno cose che non riescono, problemi che non si risolvono, cose non
capite. È vero che i docenti sono facilmente raggiungibili su Discord, ma
quella è \textbf{l'ultima chance}: bisogna imparare a \textbf{risolvere i
problemi da soli}. Se è qualcosa trattato a lezione, si riguarda la
registrazione, si chiede ai colleghi, si guardano gli appunti e le
dispense; si chiede a Google (``non riesco a installare questa cosa''), si
chiede ``alla maledettissima intelligenza artificiale''. Solo se nulla di
tutto questo aiuta, perché il problema è particolarissimo, si va dal
docente. Bisogna sforzarsi di non avere ``la pappa pronta'': fa parte del
\emph{problem solving}. È come con un datore di lavoro: non si va dal capo a
ogni problema, prima si cerca di risolverlo. Al primo anno, precisa
Fusetto, va bene scrivere qualsiasi domanda quando si vuole (e così per gli
altri docenti, secondo le indicazioni che daranno); al secondo anno
comincerà a dire ``guarda che per questa cosa potevi anche non scrivermi'',
perché nel lavoro non si può scrivere al capo ``ogni due per tre''.

Chi ha bisogno di ancora più tempo ha l'opzione \textbf{Companion},
acquistabile anche durante l'anno accademico, che dà accesso alle lezioni
per un altro anno: è pensata per chi ha davvero bisogno di più tempo, e
richiede un processo dedicato (le lezioni vengono spostate su un altro
server), con più tempo, costo e lavoro. ``Non è un obbligo, ovviamente.''

% ---------------------------------------------------------------------
\section{Trappole}
\label{p1:sec:trappole}

Poi le \textbf{trappole}, ``trappoloni''.

\paragraph{``Mi basterà ascoltare?''} Molto spesso no.

\paragraph{``Posso fare altro nel mentre?''} ``Altro trappolone enorme, che
l'anno scorso ha fatto cadere tantissime persone'': ``tanto sono indicazioni
di massima, poi mi rivedo la lezione''. No. Andrea Morini, ``come direttore
didattico, terribile'', avverte che i docenti sanno se, durante una lezione
da remoto, qualcuno guarda una serie TV o gioca; e ci sono stati studenti
che giocavano in aula mentre il docente spiegava. ``Non è ammissibile in
nessun caso, mai'': a quel punto interviene lui, ``io sono quello cattivo''.
Massima attenzione: \textbf{nessuno è multitasking}, ``è falso''. Si segue la
lezione e basta, si spegne il cellulare. Anche perché seguire più cose
contemporaneamente \textbf{affatica}, anche se non sembra.

\paragraph{``Farò due minuti di ritardo, che sarà mai?''} C'è un ritardo
ammesso di 5--10 minuti, ma \textbf{non c'è il quarto d'ora accademico},
a meno che il docente non lo conceda. Se il docente dice ``10:10 in aula'' o
``14:10 in aula'', a seconda dell'orario, o ``questo orario, \emph{sharp}'',
bisogna essere puntuali.

\paragraph{``Sono come tutorial.''} ``Questo è un grande altro tema'': il
trappolone del ``tanto è come un tutorial di YouTube, poi me lo seguo in
registrazione''. No: rispetto a un corso su YouTube, su Udemy o a un corso di
formazione registrato, si ha un vantaggio, \textbf{si può chiedere
direttamente al docente}, che è un esperto del settore (``nel corso di Udemy,
vi posso assicurare, non sapete chi è il docente''). Anche quando una
lezione è un tutorial, ``non è davvero un tutorial'': si può chiedere di
rispiegare. E, ricordando le soft skill (niente pregiudizi, rispetto), se un
collega chiede ``prof, può rispiegare?'' e noi avevamo capito, anche se
pensiamo di annoiarci, conviene pensare ``magari mi sono perso comunque
qualcosa''. ``Non è un errore rispiegare, anzi, molto meglio rispiegare.''

% ---------------------------------------------------------------------
\section{Gli obiettivi dei corsi}
\label{p1:sec:obiettivi}

\subsection{Obiettivi primari: portfolio, problemi reali, lavoro di gruppo}

Il primo obiettivo, tutto l'anno e tutti gli anni, è \textbf{lavorare per
realizzare un portfolio}. Oltre al ``pezzo di carta'', la cosa principale per
entrare nella game industry è il portfolio: i giochi creati, i prodotti
fatti, i modelli 3D per i 3D artist, il software per i programmatori, gli
script per i narrative designer, le task per i game designer, le animazioni
per i character animator.

Attenzione, ``soprattutto programmatori e 3D artist, che di solito
sottovalutano questo punto'': \textbf{i giochi pubblicati sono l'elemento di
portfolio più importante di tutti}, più degli altri giochi fatti e ancora di
più delle task. Quando un selezionatore delle risorse umane, ``ma soprattutto
un lavoratore di un'azienda videoludica'', valuta un portfolio, considera più
importanti i giochi pubblicati su itch.io o su Steam, anche con gli amici
per una jam (non devono essere tripla A, devono essere accessibili), di
qualsiasi elemento legato a giochi mai pubblicati. ``Ho fatto un modello 3D.
È stato pubblicato in un gioco? No? Vale di meno'', e vale ancora meno della
task. Per questo, andando avanti negli anni, bisognerà riuscire sempre di più
a realizzare giochi, e giochi pubblicati; e qui AIV aiuta con i Dev Projects.

Il secondo obiettivo è lavorare su \textbf{problemi e progetti reali}: c'è
``un implicito accordo'' fra docenti e studenti per cui i docenti danno
problemi che davvero si troveranno nell'industria, quelli che si affrontano
tutti i giorni. E il terzo è lavorare \textbf{in gruppo} per potenziare le
soft skill. ``Fondamentale'': bisogna creare fin dai primi mesi un
\textbf{ottimo corpo classe}, una buona sinergia fra studenti. Secondo i
numeri e le analisi di AIV (``e Andrea ve lo può confermare''), le classi
molto unite, con studenti che si supportano, hanno molte più possibilità e
meno problemi di quelle che litigano. ``È banale da dire'', eppure non si sa
quante volte micro faide e micro litigi nati per niente creino problemi
``atroci, enormi''.

Andrea Morini aggiunge che queste cose derivano spesso da una sbagliata
considerazione del proprio \textbf{ego}: ``qui non c'è ego''. Non vuole
vedere nessuno che, per motivi di ego, rallenti o danneggi i colleghi,
anche involontariamente, senza accorgersene e senza malizia, ``ma succede''.
``Qui l'ego lo potete mettere da parte, e allora andiamo tutti d'accordo.''

\subsection{Obiettivi secondari: competenze a T rovesciata}

Fra gli obiettivi secondari ci sono le \textbf{competenze a T
rovesciata}\footnote{La trascrizione dice ``competenze ti è rovesciata'' e
``competenza tirovesciata''. Il riferimento è al modello delle competenze
``a T'' (\emph{T-shaped skills}): una competenza profonda in una
disciplina (l'asta verticale) unita a una conoscenza più ampia delle
discipline vicine (il tratto orizzontale). Il relatore lo chiama ``a T
rovesciata'', con il tratto orizzontale alla base.}. Nell'industria si torna
a parlare del fatto che stanno tornando lentamente i \textbf{generalisti}.
È un ciclo costante dagli anni Settanta, da quando sono nati l'informatica e
il software: le industrie tecnologiche preferiscono, alternativamente,
persone con competenze variegate che sanno lavorare su tutto un prodotto (i
cosiddetti \emph{full stack}: programmatore full stack, game designer full
stack) oppure persone molto specializzate, \textbf{verticali}. Fino a prima
dell'IA era molto importante essere ``terribilmente verticali''; proprio in
questi giorni Fusetto e Morini si sono confrontati sul fatto che molte
aziende cominciano a pensare di tornare indietro e avere generalisti,
perché con l'IA c'è ``questa sorta di verità e illusione'' (``lo sapremo
solo con calma nel tempo'') che una persona con competenze differenti
funzioni di più.

AIV insegna a T rovesciata, e questo mette gli studenti in una sorta di
``bunker protettivo'': i corsi sono verticali, ma dentro la propria
verticalità si vedono tante cose. In game design si è verticali sul game
design, ma si vede Unreal, si programma con i Blueprint, si fa level design,
combat design: ``vedete tutto quello che c'è nel game design''. E nessuno,
nell'industria, chiederà a un game designer di fare modelli 3D: magari
competenze di base per mettere gli asset nell'engine, ``ma non vi potrete
sostituire a un modellatore 3D''. Bisogna quindi essere
\textbf{intraprendenti}: imparare tutte le sfaccettature del proprio lavoro
e, al secondo e terzo anno, nei lavori di team, imparare anche \textbf{per
osmosi} da ciò che fanno i colleghi.

Andrea Morini aggiunge un dettaglio: AIV ``si mette sempre in discussione''.
Un paio di volte l'anno va a vedere ``non so quante'' \emph{job
description}, le posizioni aperte, per verificare quanto sia vero che il
generalista sta tornando di moda: ``poco, poco''. Oggi le posizioni aperte
sono ancora ``estremamente verticali, non se ne esce''; nei piccolissimi
studi di tre o quattro persone un generalista può essere più apprezzato,
altrimenti no. Quindi, all'uscita, ci si propone come \textbf{specialisti
della propria disciplina}, e in azienda scopriranno che si sa fare anche
altro. Spesso si entra con una qualifica e si esce con un'altra (magari si
entra in un ruolo junior di base, ``adesso poi dovremo vedere che cosa
succederà con le posizioni junior''), e nulla vieta di fare carriera
all'interno.

\subsection{Competenze spendibili fuori dal settore}

Le competenze sono importanti anche per l'\textbf{extrasettore}: in tutti i
corsi si imparano cose spendibili anche fuori dalla game industry, e questo
serve per due motivi. Se non si trova lavoro nell'industria, si possono usare
altrove; e se dopo tre anni ci si convince che quello non è il proprio
futuro, idem. Fusetto ha ``moltissimi studenti'' che hanno cambiato carriera
usando le competenze acquisite. Si augura che tutti facciano il lavoro per cui
si sono iscritti, ma ``non abbiamo la palla di cristallo'': meglio prevenire
che curare.

Al terzo anno, inoltre, si fanno molte lezioni, anche con i \textbf{Career
Days}, sulle \textbf{competenze imprenditoriali} e sul
\textbf{freelancing}: non c'è solo l'assunzione in azienda, c'è il freelance
(Fusetto è freelancer) e c'è chi vuole aprire un'azienda. L'Italia, dice,
ha due cose ``fenomenali'': la capacità degli italiani di essere imprenditori
e di sognare in grande, e ``la stragrande capacità del governo italiano di
mettere burocrazia per impedire l'imprenditorialità''. Sulla burocrazia non
si può lavorare, ma si può diventare intraprendenti, evitare errori enormi e
diventare imprenditori migliori, ``per esempio degli americani'': un
americano apre una società in due ore con una carta di credito e al primo
problema si blocca; un italiano dice ``che sarà mai, ho passato ore
chilometriche a fare burocrazia''.

% ---------------------------------------------------------------------
\section{Attestato di partecipazione e certificato di eccellenza}
\label{p1:sec:eccellenza}

La parte sui certificati la presenta Andrea Morini. Da qualche tempo AIV ha
introdotto l'\textbf{attestato di partecipazione}, che attesta che si è
partecipato al corso, cioè che si è fatto ciò che il corso richiede:
\textbf{consegnare le task}. Non basta essere stati presenti (``è implicito
che è importante''); è ancora più importante \textbf{produrre}, così da potersi
confrontare su ciò che si produce.

Accanto c'è il \textbf{certificato di eccellenza}\footnote{La trascrizione
dice ``il CTV Audi Eccellenza'' e ``il certificato di scelta'': si tratta in
entrambi i casi del certificato di eccellenza.}, per chi dà importanza a
tutto ciò che si sta dicendo e applica queste best practice. I requisiti
sono tre:
\begin{enumerate}
  \item \textbf{aver prodotto molto e bene}: essere promossi, aver
        consegnato \textbf{tutte} le task dell'anno (``per alcuni corsi sono
        tante'') e superato gli eventuali test; oppure aver consegnato
        \textbf{almeno l'80\%} delle task, ma con \textbf{metà di questo 80\%
        valutata eccellente}, a discrezione del docente;
  \item \textbf{aver partecipato ai Dev Projects}, collaborando fino al
        \textbf{completamento} del progetto (non necessariamente alla
        pubblicazione); se per qualsiasi motivo non è stato possibile
        (``non perché lavorate, non potete'', ma perché il progetto è stato
        chiuso perché non andava bene, ``mille motivi''), aver partecipato
        almeno a delle \textbf{game jam}: ``è il \emph{second best}'';
  \item aver partecipato \textbf{almeno a un secondo evento di settore}
        ogni anno, oltre alla jam: ``sembra che dobbiamo metterla come un
        premio, e invece dovrebbe essere la prassi''.
\end{enumerate}

Gli eventi di settore sono fondamentali soprattutto per il
\textbf{networking}: crearsi una rete di conoscenze nell'industria che, in
futuro, ``probabilmente mi agevolerà'' nel trovare lavoro. Va fatto da
subito, non alla fine del corso: le game jam sono un modo, gli eventi di
settore (AIV ne indicherà alcuni; sulla slide ce ne sono degli esempi) sono un
altro. ``Dovete farlo''; AIV lo riconosce con il certificato. ``Dovreste
farle a prescindere, però se le fate a noi piace, tra virgolette, premiarvi
con un bollino di eccellenza.''

Il certificato ha un valore anche nel curriculum. Come Morini ha spiegato
``anche ieri all'open day'', quando si selezionano persone per un lavoro si
fa una prima scrematura cercando almeno un certificato di scuola, università
o accademia equiparata (``non importa laurea o non laurea, all'estero''):
l'importante è aver fatto un percorso dopo le superiori. Poi si guarda il
portfolio, ma si guardano anche i \textbf{meriti}. Il certificato di
eccellenza è ``un'ottima occasione per brillare''.

% ---------------------------------------------------------------------
\section{I Dev Projects}
\label{p1:sec:devprojects}

``Cerchiamo di correre verso la fine, abbiamo ancora un quarto d'ora
circa.'' ``Voi siete qua perché volete fare giochi, giusto? Quindi
facciamoli.''

\subsection{Che cosa sono}

Come si fa un gioco? Di solito ci sono tre fasi, mutuate dal cinema:
\textbf{pre-produzione}, \textbf{produzione}, \textbf{post-produzione}. Da
questo metodo AIV prende il modello di sviluppo del secondo anno.
All'inizio del secondo anno, dopo i \emph{pressure test} e i test di inizio
anno, quindi intorno a novembre (quest'anno il \textbf{18 novembre}), gli
studenti del secondo anno vengono divisi in \textbf{gruppi
interdisciplinari}, che devono fare il \textbf{pitch} di un gioco da
sviluppare nei \textbf{due anni} successivi. I migliori vengono pubblicati.

Si chiamano \textbf{Dev Projects} e sono ``sviluppo ludico reale in contesto
didattico'': a AIV non piace chiamarli progetti studenteschi, ``perché non
lo sono''. Sono supervisionati da un comitato di esperti (``una
controlunca'' nella trascrizione), di cui fa parte anche Fusetto, dal
concept alla release. I migliori vengono pubblicati ``live'' su
\textbf{Steam} e su \textbf{Epic Games Store}: dall'anno scorso, dice il
relatore, AIV è diventata una delle tre accademie in Europa, e l'unica in
Italia, a poter pubblicare i giochi degli studenti su Epic\footnote{È
un'affermazione del relatore, non verificata.}.

``Pubblicati'' non vuol dire messi online per amici e parenti: sono
prodotti \textbf{gratuiti} che vengono scaricati \textbf{fra le 60.000 e le
100.000 volte}; \emph{Steam Rush} è stato giocato da 100.000 persone; le
recensioni medie vanno da ``positive'' a ``molto positive''. ``Non stiamo
parlando di piccoli giochi che fate nel giardino di casa'', ma di prodotti
che possono fare la differenza fra essere assunti o no: molti degli autori
di \emph{Steam Rush} sono stati assunti proprio perché l'avevano nel
portfolio.

\subsection{I giochi pubblicati}

Una slide mostra i giochi pubblicati finora (``da aggiornare'').
L'elenco, nella trascrizione, è in buona parte storpiato\footnote{La
trascrizione riporta: ``l'Eder, il Beautiful, il Duality Shift \ldots\ Tito
Rebello e Duality Shift sono già stati pubblicati. Powerful viene
pubblicato a breve. Solarpunk e D'Angelo Augusto avranno una storia un po'
più lunga. Operation Remote, Bolbuletti e Steamrash sono già fuori. Gunpire
e Ancora For Atlanti \ldots''. ``Tito Rebello'' è con ogni probabilità
\emph{Beat of Rebellion}, il gioco mostrato subito dopo; gli altri titoli
non sono ricostruibili con certezza e sono riportati come appaiono.}: fra i
titoli riconoscibili ci sono \emph{Beat of Rebellion}, \emph{Duality
Shift}, \emph{Operation Remote} e \emph{Steam Rush}, già usciti; un altro
uscirà a breve; due progetti ``avranno una storia un po' più lunga''. Due
giochi, \emph{Gunpire} e un altro, sono stati anche \textbf{finalisti al Game
Development World Championship 2025}.

\subsection{Perché sono importanti}

Perché parlarne già adesso, un anno prima? Perché i Dev Projects sono ``la
vostra chance principale per lavorare in questa industria''. ``Voi siete qua
per fare giochi, dimostratelo'': dimostrare di saper lavorare in team, di
superare problemi interpersonali, di capire come lavorano gli altri
dipartimenti, di costruire un portfolio, di lavorare in modo simulativo
del mondo del lavoro.

Alla fine del terzo anno, insieme al certificato di partecipazione, AIV
consegna \textbf{la custodia del gioco} con dentro la scheda SSD o la chiave
USB con la \emph{build} (``non c'è la slide, ma la dovrei mettere''): per far
capire di aver fatto una scuola che ha fatto realizzare un prodotto reale. I
giochi vengono poi pubblicati \textbf{nei 6--8 mesi successivi} al terzo
anno, durante i quali AIV continua a seguire gli studenti
\textbf{gratuitamente}: Fusetto e \textbf{Giacomo Caneschi}, direttore di 3D
art (capitolo~\ref{p6:cap:3d}), seguono i prodotti e lo sviluppo ``perché ci
piace che i nostri studenti pubblichino''.

% ---------------------------------------------------------------------
\section{La pipeline di sviluppo}
\label{p1:sec:pipeline}

Come avviene lo sviluppo? Saperlo serve già al primo anno, perché molti
faranno progetti nella propria classe. La pipeline mostrata è questa
(figura~\ref{p1:fig:pipeline}):
\begin{description}
  \item[Pre-produzione.] Un \textbf{concept}, un \textbf{vision document},
        una fase di \textbf{ricerca}.
  \item[Vertical slice.] ``Un prototipo prototipato''\footnote{Così nella
        trascrizione: il senso è un prototipo rappresentativo del gioco
        finito, completo in una sua piccola porzione.}.
  \item[Produzione.] Design, development, prototipo.
  \item[Release candidate.] Il candidato al rilascio. Un tempo, dice il
        relatore, nell'industria si chiamava \emph{Release to Market},
        perché era la build che si metteva nei CD; oggi che molti giochi
        sono pubblicati in digitale si chiamerebbe \emph{Release to
        Marketing}, la build consegnata al marketing per la
        promozione\footnote{La sigla in uso nell'industria del software è
        RTM, \emph{Release to Manufacturing}: la versione consegnata alla
        produzione fisica dei supporti. ``Release to Market'' e ``Release to
        Marketing'' sono le espansioni date dal relatore; la seconda, in
        particolare, sembra una sua interpretazione.}.
  \item[Polishing, marketing, community management.] Chi segue, come il
        relatore, \emph{Control Resonant} può vedere quanto marketing si
        sta facendo per spingere il gioco: ``tantissimo''.
  \item[Customer care.] Un gioco è un software, e serve assistenza al
        cliente: magari un cliente l'ha comprato ma non riesce a
        scaricarlo.
\end{description}
Un gioco, insomma, non finisce al rilascio: ci sono i DLC, ci sono le patch.
``Il \emph{Control} è già stato patchato tre volte, siamo alla 1.4, è uscito
da una settimana''\footnote{La trascrizione dice ``il controverso è già
stato patchato'': dal contesto è \emph{Control Resonant}, citato poco prima.
I dettagli sulle patch sono quelli riferiti dal relatore.}. ``C'è da
lavorare anche dopo.''

Una precisazione terminologica, valida per tutto l'anno: ``io parlo in
termini tecnici, se non è chiaro, chiedete''. Per \textbf{build} si intende
l'eseguibile del gioco, il file \texttt{.exe} che si può cliccare e lanciare
(per esempio da Steam); più in generale, la build è \textbf{la versione del
gioco su cui si sta lavorando}.

\begin{figure}[htbp]
\centering
\begin{tikzpicture}[
  fase/.style={draw, rounded corners, align=center, font=\footnotesize,
               minimum height=1.25cm, text width=2.05cm},
  >=Stealth, node distance=0.35cm]
  \node[fase] (pre) {\textbf{Pre-produzione}\\ concept, vision,\\ ricerca};
  \node[fase, right=of pre] (vs) {\textbf{Vertical\\ slice}};
  \node[fase, right=of vs] (prod) {\textbf{Produzione}\\ design, sviluppo,\\ prototipo};
  \node[fase, right=of prod] (rc) {\textbf{Release\\ candidate}};
  \node[fase, right=of rc] (post) {\textbf{Post-produzione}\\ polishing, marketing,\\ community, assistenza};
  \draw[->, thick] (pre) -- (vs);
  \draw[->, thick] (vs) -- (prod);
  \draw[->, thick] (prod) -- (rc);
  \draw[->, thick] (rc) -- (post);
  \draw[->, thick, dashed] (post.south) -- ++(0,-0.6) -| node[pos=0.25, below, font=\scriptsize]
       {patch, DLC: il lavoro continua dopo il rilascio} (prod.south);
\end{tikzpicture}
\caption{La pipeline di sviluppo dei Dev Projects come la presenta il
relatore. Ricostruzione schematica.}
\label{p1:fig:pipeline}
\end{figure}

% ---------------------------------------------------------------------
\section{Un caso: l'evoluzione di \emph{Beat of Rebellion}}
\label{p1:sec:beat}

Per far vedere come cambia un gioco lungo la pipeline, il relatore mostra
\emph{Beat of Rebellion}, uno dei giochi pubblicati l'anno precedente.
\begin{itemize}
  \item \textbf{Pre-produzione.} I bozzetti, il \emph{concept document} e
        una ``build ipotetica'', il prototipo realizzato nella fase di
        ricerca\footnote{La trascrizione dice ``la GIM che avevano nella fase
        di ricerca'': la parola non è ricostruibile; dal contesto si tratta
        della build prototipale di quella fase.}.
  \item \textbf{Vertical slice.} Dopo circa sei mesi, da gennaio a giugno,
        la vertical slice. La slide mostra il \emph{level design} visto
        dall'alto: ``non è così il gioco'', che è isometrico e segue il
        personaggio.
  \item \textbf{Produzione.} Il lavoro sul livello: il livello c'era già,
        ma è stato annotato, è stato sviluppato il personaggio e l'aspetto
        grafico. Si vede il livello crescere, e si nota la differenza a
        livello di \textbf{texturing}: un modello 3D senza texture, cioè
        senza quell'aspetto di superficie, e lo stesso modello con la
        texture. ``Queste parole imparatele tutti, non solo i grafici 3D.''
        È stato fatto anche \emph{character development} degli NPC.
  \item \textbf{Release candidate.} Ciò che hanno prodotto alla fine del
        terzo anno, quando i supervisori hanno detto ``potete pubblicarlo
        entro fine anno'': non è la versione pubblicata su Steam, ``che è
        ancora meglio'', e ha ancora qualche piccolo bug. Si vede che lo
        stanno giocando dall'editor di Unreal, perché il personaggio lascia
        delle scie, dei \emph{ghost}, legati ai test in corso. ``Ma guardate
        questa camera che fa vedere il livello'': non sono cose che si vedono
        di solito in un gioco di studenti, ``perché questo non è un gioco di
        studenti''. Fare le camere, le animazioni, capire come fare il salto,
        il doppio salto, il \emph{dash}: tutte cose che si vedranno durante
        gli anni.
  \item \textbf{Post-produzione.} Polishing, marketing, miglioramento della
        pagina del gioco ``a livello di SEO'', fatto con il relatore: per un
        periodo, ``forse anche adesso'', cercando il titolo il gioco usciva
        come primo risultato, ``c'è tutta una tecnica per farlo''. Sono
        state migliorate anche le \textbf{icone} dei poteri dei personaggi,
        che prima erano ``semplicemente abbozzate''.
\end{itemize}
Il video è senza audio (``no, non funziona l'audio''), ``ed è un peccato,
perché è un gioco che si basa sul beat'', sulla musica. ``Cercatelo online,
giocatelo'': è migliorato tantissimo.

% ---------------------------------------------------------------------
\section{Perché parlarne oggi}
\label{p1:sec:perche}

Perché dirlo il primo giorno? Perché si è qui per fare videogiochi, che sia
entrando nell'industria, come sviluppatori indipendenti o come freelance:
se si vogliono fare giochi, i Dev Projects ``sono il vostro futuro''. Di
più: al secondo e al terzo anno Fusetto obbligherà a fare le \textbf{jam},
come la \emph{Fast Starter Jam} di inizio anno\footnote{Nel
capitolo~\ref{p2:cap:gamedesign} la stessa jam è chiamata ``Firestarter''; il
nome esatto non è ricostruibile con certezza.}, ``perché voi dovete fare
giochi'': si deve uscire da AIV con \textbf{almeno due o tre giochi
pubblicati} su itch.io o su Steam.

C'è un anno prima di cominciare i Dev Projects, e in quell'anno si può
perdere entusiasmo. Il relatore vuole far capire \textbf{perché} lo si fa,
così da arrivare al secondo anno con l'entusiasmo di oggi: chi oggi non ha
l'\emph{hype} del ``io voglio farlo, quel videogioco'' non sopravviverà al
secondo anno. Bisogna mantenere l'entusiasmo: ``occhi sulla meta, come si
dice in America''.

Che cosa fare quest'anno per prepararsi al meglio?
\begin{itemize}
  \item \textbf{comprendere meglio la propria disciplina}: il primo anno si
        imparano le basi;
  \item essere \textbf{aperti a punti di vista} diversi;
  \item lavorare sul \textbf{mindset}, soprattutto su come prendere e
        recepire i \textbf{feedback}: introspettarli, farli propri, usarli,
        capirli;
  \item cominciare a \textbf{stringere amicizie} nel settore partecipando
        agli eventi, che di solito vengono comunicati sul Discord di AIV
        (ci sono chat dedicate; spesso AIV stessa viene invitata come
        accademia); ma anche informarsi da soli, scrivere, andare online:
        ``non siate passivi, siate proattivi''.
\end{itemize}

% ---------------------------------------------------------------------
\section{Falsi miti e cose a cui stare attenti}
\label{p1:sec:miti}

Le ultime slide, ``rapide'', elencano alcuni falsi miti.
\begin{description}
  \item[``Non mi serve il Dev Project, ho un portfolio da scrittore o da 3D
        artist.''] \textbf{Il gioco pubblicato batte tutto.} Nazareno
        Urbano, amico del relatore e \emph{principal environment artist} in
        Remedy, ha detto in più interviste: ``la prima cosa che guardo in un
        3D artist è quanti giochi ha pubblicato e quali. Poi il portfolio,
        ma prima guardo i giochi''.
  \item[``Posso ghostare gli altri membri del team, tanto chi se ne
        frega.''] No: si rispetta l'impegno degli altri. Se ci si dà un
        appuntamento con i compagni a una certa ora, ci si va: ``questa non
        è una scuola superiore, questa è un'università''. Si comunica fra di
        sé.
  \item[``Posso operare in autonomia.''] No: si coinvolge il team.
  \item[``Facciamo un open world MMO online multiplayer.''] No: ``fate il
        passo lungo quanto la gamba''. Niente prodotti troppo grandi: ``non
        siete qua per fare il nuovo Pokémon''; per quello servono soldi,
        tempo e persone. Bisogna pensare giochi \textbf{realizzabili in due
        anni}. Vale anche al primo anno, ``perché anche quest'anno farete
        cose''.
  \item[``Non serve la documentazione, non serve tenere traccia delle
        decisioni.''] Vale anche per chi fa progetti al primo anno: ``tanto
        non serve scrivermi come ho fatto''. Serve, ``perché se no fate un
        buco nell'acqua ogni volta''. La documentazione si fa.
  \item[``Lo facciamo quando abbiamo tempo.''] No: bisogna
        \textbf{programmarsi la giornata}, ``programmatevi la vostra
        vita''. La gestione del tempo e delle task è fondamentale e fa
        parte delle soft skill: tutte quelle capacità non tecniche di
        gestione, di \textbf{autogestione}, di rapporti interpersonali.
        ``Anche lavorare su se stessi è fondamentale.''
\end{description}

% ---------------------------------------------------------------------
\section{Chiusura e domande}
\label{p1:sec:domande}

L'ultima slide invita a sfruttare gli \textbf{eventi Growth}\footnote{È il
nome, nella trascrizione, della serie di incontri di inizio anno di cui fanno
parte i Dev Days.} per capire come lavoreranno i colleghi e ragionare sulle
altre discipline: partecipare, o almeno vedere le registrazioni. Oggi sono
collegate \textbf{31 persone su 100}: ``è un terzo, è pochissimo''; l'anno
scorso erano almeno la metà. Il relatore si aspetta di vederne almeno una
cinquantina il giorno dopo, alla presentazione di narrative design. Le slide
verranno mandate il prima possibile, caricate sulla piattaforma o inviate
in altro modo, a volte in PDF su Discord.

\subsection{RAM: 16 o 32 GB?}

Una domanda ``off topic'': uno studente ha un PC con 16 GB di RAM e uno slot
libero; conviene montare un altro banco da 16 per arrivare a 32? ``Che corso
fai?'' Game design. ``Direi di sì.'' Più o meno tutti i corsi usano Unreal
Engine, non nella stessa misura, e Unreal è ``molto pretenzioso, anche
troppo''. Il consiglio è di non stare sotto i 16 GB e di andare sui 32,
``senza fretta, però'', nell'arco dei prossimi mesi; ``c'è Natale'',
scherzano i relatori, e un banco di RAM ``oggigiorno è un ottimo regalo'' (non
è semplice trovarle: ``ormai sono firmate Morellato, le RAM''). Fusetto
aggiunge che 16 GB ancora reggono, ma 32 sono il minimo; lui, per portarsi
avanti, ne ha 64, ``non è obbligatorio, però non si sa mai''.

\subsection{Il colloquio personale e la modalità playtest}

``Potete anche usare il microfono, non è vietato, non abbiate paura.'' Il
\textbf{colloquio personale} si può usare durante l'anno: bisogna scrivere
in segreteria, che metterà in contatto con il docente più adatto al proprio
corso. Il consiglio è di usarlo \textbf{entro gennaio}, perché di solito è
il momento in cui le persone cominciano a capire se il corso fa per loro, o
addirittura se vogliono cambiare corso.

Per chi si è iscritto con la \textbf{modalità playtest}\footnote{Dal
contesto, un'opzione di iscrizione che permette di provare più corsi nelle
prime settimane e di scegliere poi quale tenere (``accendere'') e quale
lasciare (``spegnere'').}, il momento per decidere quale corso tenere è
\textbf{intorno al 30 ottobre}, perché le lezioni cominciano ufficialmente
il \textbf{19 ottobre} (narrative comincia prima, ``questo venerdì: mi
spiace, soffrirete un po' di più degli altri''). È meglio decidere almeno
sette o otto giorni dopo l'inizio delle lezioni, ``tendenzialmente lo
capisci''. Ma se, per esempio, si fa il playtest di game design e
programmazione, e alla prima lezione la matematica ``non va giù'' e la
programmazione non fa per sé, non serve aspettare fine ottobre: il colloquio
può servire anche a questo.

\subsection{Il resto della settimana}

L'appuntamento successivo è \textbf{alle 11:30} di oggi, con l'incontro di
game design (capitolo~\ref{p2:cap:gamedesign}); domani la presentazione di
narrative design con Fusetto (capitolo~\ref{p3:cap:narrative}); giovedì la
comprensione del testo (capitolo~\ref{pt:cap:comprensione}). ``Che sia l'inizio buono di un percorso fichissimo.
Buon inizio settimana, forza e coraggio!''
